\section{Discussion, Limitations \& Future Work}
\label{discussion}

% ── 1. Interpretation ─────────────────────────────────────────────────
%   LLM advantage (if observed): biological priors (collapse avoidance,
%   expected footprint, motif identity) give the LLM a head start that
%   TPE must learn from data; delta-driven reversal of bad moves provides
%   implicit gradient information.
%   Optuna advantage (if observed): 7-parameter space is small enough for
%   TPE to map quickly; LLM's one-param-at-a-time conservatism may limit
%   multi-dimensional exploration; budget asymmetry (12–15 vs 10–12 iters)
%   may confound the comparison.
%   Broader point: LLM advantages likely scale with problem complexity.

% ── 2. Biological interpretation ──────────────────────────────────────
%   • RBFOX2_K562 (hardest): wider bandwidth (>50 nt) helped — suggests
%     diffuse crosslink signal or broad footprints.
%   • Cell-line effects: if optimal parameters differ between RBFOX2_K562
%     and RBFOX2_HepG2, this validates the dataset-dependent tuning premise.
%   • use_input_covariate: note whether enabling it consistently improved
%     or hurt different datasets.

% ── 3. Limitations ────────────────────────────────────────────────────
%   (a) chr21-only: ~1.5% of the genome; results may not generalise to
%       full-genome landscapes with different gene density and repeat content.
%   (b) Small search space (~10^5 combinations): limits how much LLM
%       reasoning can outperform systematic sampling.
%   (c) Single LLM, no prompt ablation: only DeepSeek tested; no ablation
%       of priors, delta table, or temperature.
%   (d) Budget asymmetry: LLM 10–12 iters vs. Optuna 12–15 — not iso-budget.
%   (e) Motif background: mononucleotide shuffle destroys dinucleotide
%       context; dinucleotide shuffling would be more rigorous.
%   (f) Benchmark reference: ENCODE "top regions" are themselves pipeline
%       output, not manually curated gold standard.

% ── 4. Future work ────────────────────────────────────────────────────
%   (a) Genome-wide validation of best chr21 parameters.
%   (b) IDR-based reproducibility (ENCODE standard) to replace simple
%       overlap.
%   (c) Dinucleotide-shuffled motif background.
%   (d) LLM ablation: remove priors / delta table; test other models;
%       iso-budget re-run.
